\chapter{常見問題與處理方法}\label{ch:troubleshooting}

\chapterlead{遇到問題時，先認清它發生在哪一步：進入應用、讀取檔案、編輯模型、執行計算，還是儲存結果。找到這一步，再做一次小而明確的檢查，通常比同時修改許多設定更有效。本章按日常使用中看到的現象組織，具體操作可返回相應教程。}

\section{安裝後無法啟動}
\index{Anzhuang yu qidong@安裝與啟動問題}

先核對安裝包的作業系統和處理器架構，再看安裝器要求的 MATLAB\regmark{} Runtime 發行版。應用版本與Runtime發行版分別編號，獨立應用需要匹配的Runtime環境。

如果安裝器仍在配置 Runtime，應等它明確顯示完成。Windows 離線包提示缺檔案時，把 EXE 與同名分卷放在同一目錄；macOS 僅應用包沒有檢測到環境時，在安裝嚮導中選擇正確的 Runtime 根目錄。安裝後仍失敗，保留啟動提示和安裝位置，使用第\ref{ch:remote}章介紹的幫助入口。

\section{點選命令後，好像沒有反應}

先檢視是否有尚未完成的引數視窗、確認視窗或系統檔案選擇視窗。模態視窗開啟時，主視窗暫時不能接受其他操作。尤其在 macOS 上，系統檔案面板可能屬於另一個執行環境視窗；切換應用視窗或檢視 Dock 中的視窗，找到檔案選擇面板並完成或取消，便可繼續。

執行計算或載入大結構時，則檢視\ui{模擬}面板或狀態區域。如果訊息仍在更新，說明任務還在進行。任務進行時等待狀態更新，完成後再執行下一項操作。

\section{線上版一直停在連線頁面}
\index{Zaixian ban lianjie@線上版連線}

回到固定入口，檢視工作空間是否還在啟動。第一次啟動可能較慢，保持一個活動標籤頁等待即可；已經在另一個標籤頁開啟的工作空間，優先回到原標籤頁繼續。入口提供\ui{重試連線}時，先使用它，而不是連續重新整理多個網頁。

若平臺顯示名額已滿，按入口提示稍後重試。確實需要\ui{重啟工作空間}時，先儲存當前專案；重啟會結束記憶體中的會話，已儲存雲端檔案仍保留。網路中斷時使用重連；只有確認需要結束當前例項時才重啟。

\section{上傳了檔案，專案中卻沒有結構}

線上版的\ui{檔案}面板與專案瀏覽器管理不同內容。上傳只是把本機檔案放入雲端目錄，還需要雙擊該檔案或使用\ui{匯入}，才能把結構加入當前專案。先確認上傳已完成，再觀察專案瀏覽器的 \ui{Structure} 分支。

若匯入後仍沒有結構，檢查檔案內容與所選格式是否一致。修改副檔名不會把一種格式轉換成另一種格式。先用隨書的小型 \mbox{\file{silicon.cif}} 或 \mbox{\file{hydrogen.xyz}} 檢查匯入入口，再處理自己的檔案；如果示例能開啟，重點檢查原檔案的晶胞、元素、座標與格式說明。

\section{螢幕上的原子數與預期不同}
\index{Zhouqi yingxiang@週期映像與原子數量}

三維檢視中的 \ui{sites} 是模型獨立位點數，\ui{atoms} 可包含邊界和週期映像。核對超胞和缺陷時，使用獨立位點數。例如 8 位點晶胞建立 $2\times2\times2$ 超胞後應有 64 位點，而顯示原子球可以更多。

若位點數本身不符，檢視是否只改了 \ui{Repeat cell} 顯示重複，或者是否對已經擴大的模型又執行一次超胞。再檢查檔案使用的是原胞還是常規晶胞。對於按空間群構建的結構，非對稱位點還會展開為等價位置，輸入一行並不表示輸出只能有一個原子。

\section{找不到建模命令，或命令不可用}

先啟用三維結構文件。\ui{晶體建模}和\ui{分子建模}根據當前文件出現，切到工作流或圖形後，相應頁面可能不再顯示。然後檢查物件型別與選擇數量：距離需要兩個原子，鍵角需要三個，二面角需要四個，精確移動一個原子則應只有一個位點。

選擇不清楚時，按 Esc 退出當前工具，清除選擇後重新按順序點選。週期映像較多時，可先把顯示重複恢復為 $1\times1\times1$。如果引數框已開啟，先核對框中的位點序號，再決定是否重新選擇。

\section{預覽正確，關閉視窗後修改卻消失了}

預覽只用於檢查，\ui{取消}或關閉對話方塊會恢復此前文件。要接受修改，應點選\ui{應用}；要在重新開啟專案後保留，還要儲存專案。這三個階段見第\ref{ch:modeling}章。

應用後可以撤銷，儲存後則以最近儲存狀態為新基準，當前撤銷歷史清空。若需要長期保留前後對照，在重要修改前另存專案或保留獨立結構。預覽過期通常表示文件在預覽建立後改變，重新開啟命令並預覽即可，不必反覆應用舊視窗。

\section{建模命令報告錯誤}

表面構建先核對 Miller 指數不是全零，厚度與真空為合適正值，終止候選序號在可用範圍內。正交化失敗時，可能是當前長度、角度和原子數限制下沒有合適整數超胞，可有目的地調整搜尋範圍。

若提示明確指出某個已安裝元件無法載入，例如“無效的 MEX 檔案”，這是應用元件的載入問題。保留輸入和完整提示，可先匯入已經準備好的表面結構繼續學習，再向支援人員說明版本和平臺。

裝填失敗則優先檢查空間條件：降低初始密度或分子數、擴大允許區域，或更換種子。生成後仍有未解決的近距離接觸，應先處理這些接觸再進入後續平衡，而不是隻增加計算步數。

\section{工作流執行了，但沒有得到曲線}

先看當前活動工作流是否正確，再檢查曲線節點是否連線到相應計算輸出。\ui{Energy Curve} 是結果繪圖步驟，不會獨立完成 SCF。計算節點之前的結構輸入或引數若失敗，後面的曲線自然沒有可用資料。

展開 \ui{Results} 中的 \ui{Datasets} 與 \ui{Plots}，確認結果是否已生成而文件尚未開啟。閱讀\ui{模擬}面板末尾的提示，找出最後完成的節點。必要時先回到最簡單的“結構輸入—SCF—曲線”流程，確認資料能順利傳遞，再加入其他分析。

\section{\headinglatin{SCF} 達到最大步數仍未收斂}
\index{Shoulian panduan@收斂判斷}

先區分兩種趨勢：誤差持續下降但尚未到閾值，或誤差長期振盪。前者可以考慮增加 \ui{Max SCF steps}；後者則應檢查結構、佔據和混合設定。增加最大步數只允許繼續迭代，不會自動解決振盪。

\begin{steps}
\item 檢查原子是否重疊、鍵長是否異常、模擬盒是否明顯過小。
\item 核對元素、贗勢、分子或晶體型別，以及截斷能的單位。\ui{Ecut (Hartree)} 中的 20 表示 20 Ha。
\item 對金屬或接近簡併的體系，檢查電子佔據與展寬設定是否符合任務；不要照搬孤立氫分子的零溫設定。
\item 檢視誤差和能量隨迭代的變化，每次只調整一個主要設定，保留前一次輸出供比較。
\end{steps}

能量曲線看起來平坦不等於電子密度已滿足收斂條件。先讀結束訊息，再讀最終誤差。幾何最佳化還要檢查原子力與幾何停止條件；電子收斂和幾何收斂應分別判斷。

\section{計算結果與教程不同}

先比較結構、泛函、贗勢、截斷能和 $k$ 點，再比較繪圖參考。能帶以費米能為零時，應與採用同一零點的曲線比較。不同晶胞的高對稱路徑也可能不同，應同時核對路徑座標。

本書DOS採用四面體積分，網格和能量取樣會影響曲線細節，額外能帶影響能夠觀察的空態範圍。最佳化結果還與初始構型、固定原子和收斂條件有關。開始練習時，先重現隨書輸入，再逐項改變引數。軟體版本和初始電子態可能改變迭代次數或最後幾位數值，檢查物理量和收斂程度比要求日誌逐行完全相同更有意義。

\section{體資料為空白，或者顏色不明顯}

先核對當前通道和資料範圍。等值面閾值超出範圍時不會形成表面；只含正值的通道也不會憑空產生負等值面。回到 \ui{Value Distribution} 選擇合適閾值，確認構建完成。

切面顏色接近時，適當縮小顏色範圍，並檢查 Miller 指數；I、J、K 共同定義一箇中央平面，不能全為 0。裝置無法支援直接體渲染時，按狀態提示使用切面。插值和平滑改善外觀，不能增加源網格本來沒有的細節。

\section{儲存了專案，卻在本機找不到}

桌面應用儲存到所選本機目錄；線上版儲存到雲端目錄，還需要下載才得到本機副本。已經儲存過的專案繼續寫回原位置，另選資料夾不會自動改變儲存目標；改名或換目錄應使用\ui{另存為}。

線上版下載前先儲存，再在\ui{檔案}面板選擇對應 \file{.ks}。檢視瀏覽器下載列表，確認傳輸和本機儲存完成。單個檔案和一次傳輸的大小限制見第\ref{ch:online}章。登出前確認重要內容已寫入檔案，不依賴記憶體會話長期保留。

\section{遠端作業完成，卻沒有匯入結果}

本節面向完整 MATLAB\regmark{} 工具箱使用者；獨立應用與基於 Runtime 的線上版不提供同一遠端選單。開啟原提交專案，在\ui{遠端作業與結果}中重新整理狀態，檢視 Diary，再執行\ui{取回並匯入}。原工作流應仍在專案中，便於結果對應。

Unknown 表示暫時無法明確讀取狀態，不等於失敗。\ui{刪除記錄}也不等於取消作業；需要結束計算時使用\ui{取消作業}。取回後檢查計算收斂狀態，並儲存專案。

\section{對話問題被當成命令執行}

先檢查輸入開頭。命令視窗將普通文字按命令解釋；需要提問時，在自然語言前加一個美元符號，獨立問題可用兩個美元符號。具體示例見第\ref{ch:remote}章\ref{sec:chat-use}節。若輸入方式正確但提示缺少配置，再開啟設定檢查服務、地址和模型。

\section{新增任務沒有得到預期資料}

先核對任務的作用。NSCF節點準備選項，仍需後續Band Calculation或DOS執行取樣；布里淵區繪圖只顯示倒空間幾何，不計算帶隙；TDDFT需要基態空軌道，圖形節點的輸入處理限制及完整MATLAB\regmark{}中的求解方法見第\ref{sec:tddft-workflow}節。對照第\ref{ch:workflows}章表\ref{tab:task-catalog}檢查完整流程，保留出錯前最後一條日誌。檢視結果時，核對圖形所屬的執行記錄。


\section{準備一份容易理解的問題說明}

寫明軟體版本、作業系統、使用方式、輸入檔案和最短操作步驟，再寫預期與實際結果。附上問題發生時的完整提示、相關日誌和一張保留引數的截圖；能夠重現的小專案通常比大量零散截圖更有幫助。

從\ui{幫助}開啟日誌，按發生時間找到相關部分。使用\ui{請求支援}會開啟郵件草稿，檢查內容後由自己傳送。問題解決後儲存最終採用的設定，並重新執行一個小例子，確認可以繼續工作。
